\section{Results from the literature}\label{sec:literature}
% Published results used by the proof.  All of them are proved inside the formal development.
\begingroup
\emergencystretch=4em
\providecommand{\literatureformin}[1]{\textup{[}\textit{#1}\textup{]}}
\providecommand{\literatureform}[1]{\unskip\nobreak\hskip0.4em plus 0.2em\literatureformin{#1}}
\providecommand{\literaturepar}[1]{\par\medskip\noindent\textit{#1.}\enspace}
\providecommand{\literaturegroup}[1]{\multicolumn{4}{@{}l}{\textit{#1}}\\[1pt]}

The proofs of Theorems~A and~B use a large body of published mathematics.  This section records
the published results on which they rely, grouped by subject.  None of them is an axiom of the
formal development or a hypothesis of the main theorems: each is stated as a proposition of the
development \cite{LeanSource2026} and proved there (\cref{sec:verification}, \cite{Sutherland2026}).
The citations identify the mathematical source of a statement; the formal proofs often follow
routes different from the published ones.

Most of these statements are definitions in the namespace \texttt{GlobalTheory.Published} (a few
in \texttt{UpperTheory.V.Published}); we write $\sfname{X}$ for the statement named \texttt{X} in
one of these namespaces, unless another place is indicated.  Of the $332$ statements so labelled, $250$
are conclusions of theorems of the development; the others are components of larger statements or
hypotheses of auxiliary results, and the prominent ones among them (for instance the full local
Langlands correspondence for $\GL_n$) are replaced in the proof by the statements described in
\cref{sec:literature:forms}.  Further classical theorems (the zero-free regions, Kottwitz's integral
models, the comparison theorems of $p$-adic Hodge theory, resolution of singularities, \dots) are
proved without this label.  The results enter the proof through two packages of propositions,
\texttt{UnitaryGroupThirdPlaceHypotheses} ($60$ fields) and \texttt{ThetaLocalInputs} ($53$
fields), from which the hypothesis package of \cref{sec:families} is built; through the fifteen
inputs of the covolume bound of \cref{sec:covolume}, one of which is a package of $28$ statements;
and inside the proofs in \cref{sec:top,sec:window,sec:galois}.

For each result we give the statement in the form used, references, and a mark:
\literatureform{same} if the formal statement is the published theorem in its stated generality,
up to normalisations and the choice of formal models of the objects involved;
\literatureform{special case} if it is restricted, for instance to the groups $\Uni(p,q)$, to
particular number fields or to a range of absolute convergence; \literatureform{variant} if it is
formulated differently, stronger in some respect, weaker in another, or packaged in another way.
A comparison with the sources classifies $246$ of the $332$ labelled statements as of the first
kind, $41$ as special cases and $41$ as variants; four have no precise published counterpart
(\cref{sec:literature:forms}).  Some of the
formalisations are long: those of the toroidal compactifications of unitary Shimura varieties, of
the integral toroidal compactifications of Siegel moduli, of Kottwitz's integral models, of Weil's
Siegel formula for unitary groups, of the arithmetic Hilbert--Samuel theorem and of the $p$-adic
monodromy theorem use about $94$, $69$, $66$, $34$, $30$ and $24$ thousand lines of Lean that are
not shared with other proofs.  \Cref{tab:literature} summarises the main results.

\subsection{Analytic number theory}\label{sec:literature:ant}

\begin{description}[leftmargin=0pt,itemsep=1pt,parsep=0pt,font=\normalfont\itshape]
\item[Dirichlet $L$-functions] Mathlib \cite{Mathlib2020,LoefflerStoll2025}, \cite{Davenport2000}.
  Continuation, change of level, the functional equation
  $\Lambda(\chi,1-s)=N^{s-1/2}\varepsilon(\chi)\Lambda(\bar\chi,s)$ for primitive $\chi$ modulo $N$,
  and $L(\chi,s)\neq0$ for $\Re s\ge1$ (Hadamard, de la Vall\'ee Poussin).  \literatureform{same}
\item[Zero-free regions] \cite{delaValleePoussin1899}, \cite[Ch.~13--14]{Davenport2000}.
  $\zeta(s)\neq0$ for $\Re s\ge1-c/\log(|\Im s|+2)$, $s\ne1$, and the prime number theorem with error
  term $O(y\,e^{-c\sqrt{\log y}})$; for each $q$, at most one zero of $\prod_{\chi\bmod q}L(s,\chi)$
  lies in $\Re s\ge1-c/\log(q(|\Im s|+2))$, and it is a simple real zero of a real character.
  \literatureform{same}
\item[Siegel's theorem and the Siegel--Walfisz theorem] \cite{Siegel1935b,Walfisz1936}.
  $|L(1,\chi)|\ge c(\varepsilon)q^{-\varepsilon}$ for real non-principal $\chi$ modulo $q$, hence
  $\beta\le1-c'(\varepsilon)q^{-\varepsilon}$ for real zeros; for every $A>0$,
  $|\psi(x,\chi)-\delta_\chi x|\le C_Ax\,e^{-c_A\sqrt{\log x}}$ for characters $\chi$ modulo
  $q\le(\log x)^A$, where $\delta_\chi=1$ if $\chi$ is principal and $0$ otherwise, and for
  $q\le e^{c\sqrt{\log x}}$ once the term of a possible exceptional zero is removed.
  \literatureform{same}
\item[Zero repulsion and mean values] The supply of host fields (\cref{thm:supply}) also uses a
  Deuring--Heilbronn type repulsion, a log-free zero-density estimate near $\Re s=1$ and a mean-value
  estimate for real character sums, of the types in \cite[Ch.~18]{IwaniecKowalski2004} and
  \cite{HeathBrown1995}, formulated as needed.  \literatureform{variant}
\item[Linnik's theorem] \cite{Linnik1944}.  The least prime $p\equiv a\bmod q$ is at most $Cq^L$
  ($\sfname{Linnik}$).  \literatureform{same}
\item[Hecke $L$-functions] \cite{Hecke1920,Tate1967thesis}.  For a unitary Hecke character $\chi$
  which is not of the form $|\cdot|^{it}$, $L(s,\chi)$ and its partial versions (finitely many
  Euler factors removed) continue to entire functions, the latter without zeros on $\Re s\ge1$, and
  $L(1,\chi)\neq0$ for $\chi\neq1$; Tate's functional equation with explicit $\varepsilon$-factors for
  characters of CM type.  \literatureform{same}
\item[Partial zeta functions] \cite{Klingen1962,Siegel1970,DeligneRibet1980}.  For totally real
  fields: the functional equation, rationality of the values at negative integers, and the
  $\ell$-adic congruences of Deligne--Ribet.  \literatureform{same; Deligne--Ribet: variant}
\item[Chowla--Selberg formula] \cite{ChowlaSelberg1967}.  For imaginary quadratic $K$ of
  discriminant $-d<-4$, the mean over the ideal classes of $\log(4\pi\Im\tau\,|\eta(\tau)|^4)$ equals
  $\gamma+\log(2\pi/\sqrt d)-L'(1,\chi_{-d})/L(1,\chi_{-d})$; it gives the periods of CM elliptic
  curves used in \cref{sec:covolume}.  \literatureform{same}
\end{description}

\subsection{Class field theory and Galois cohomology}\label{sec:literature:cft}

\begin{description}[leftmargin=0pt,itemsep=1pt,parsep=0pt,font=\normalfont\itshape]
\item[Class field theory] \cite{Artin1927,Tate1967gcft,Chevalley1951}.  The global Artin map
  $\A_E^\times\to G_E^{\mathrm{ab}}$, an open quotient map with kernel the closure of
  $E^\times(\A_{E,\infty}^\times)^\circ$ sending local units onto inertia and uniformisers to geometric
  Frobenius elements ($\sfname{GlobalArtinReciprocity}$); ray class fields; Chevalley's theorem on
  units, from which local class field theory is derived; the Hasse norm theorem for quadratic
  extensions.  \literatureform{same}
\item[Chebotarev's theorem] \cite{Chebotarev1926}.  In qualitative form: for $L/F$ finite Galois
  (or $U\subset G_F$ open normal), $g$ and a finite set $S$, some unramified place $v\notin S$ of
  degree one has Frobenius $g$ (resp.\ in $gU$) ($\sfname{ChebotarevDegreeOne}$).
  \literatureform{special case}
\item[Complex multiplication] \cite{Shimura1971,Silverman1994,SerreTate1968,Gross1980}.  Hilbert and
  ray class fields of imaginary quadratic fields through $j$ and Weber functions, Shimura's
  reciprocity law, good reduction of CM elliptic curves over $\Q_\ell^{\mathrm{ur}}$ for
  $\ell\nmid6d_K$, Gross's curves with minimal discriminant $-p^3$.  \literatureform{same}
\item[Hermitian forms, strong approximation]
  \cite{Landherr1936,Jacobowitz1962,Kneser1966,PlatonovRapinchuk1994}.  Classification and Hasse
  principle for hermitian forms over CM and $p$-adic fields; strong approximation for $\mathrm{SU}(J)$
  with respect to $S$ when $\mathrm{SU}(J)(F_S)$ is non-compact.  \literatureform{same}
\item[Local Galois cohomology] \cite{Tate1963,NeukirchSchmidtWingberg2008}.  For finite
  $\Z/\ell^r$-modules of decomposition groups of number fields: finiteness, local duality, the
  Euler characteristic formula, orthogonality of unramified classes.  \literatureform{same}
\item[Global duality] \cite{Tate1963,Milne2006,Wiles1995}.  For finite $\Z/\ell^r$-modules $M$ of
  $G_{F,S}$, $\ell$ odd or $F$ totally complex: the Poitou--Tate nine-term sequence with normalised
  invariants, $\Sha^2(M)\cong\Sha^1(M^\vee(1))^\vee$ and $H^n(G_{F,S},M)=0$ for $n\ge3$
  ($\sfname{PoitouTateNormalised}$); Tate's global Euler characteristic formula; the
  Greenberg--Wiles formula.  \literatureform{same}
\item[Ferrero--Washington theorem] \cite{FerreroWashington1979}.  For abelian $K\ni\mu_\ell$
  ($\mu_4$ if $\ell=2$) the $\ell$-rank of the class group of $K(\mu_{\ell^n})$ is bounded in $n$.
  \literatureform{same}
\item[Kato's theorem] \cite[Thm.~14.2]{Kato2004}.  Vanishing of $H^1_f(\Q,W)$ and
  $H^1_f(\Q,W\otimes\theta)$ for the twists $W$, with Hodge--Tate weights $\pm e$, of the
  representations of $\ell$-ordinary newforms of weight $2e+1$ and quadratic nebentypus $\theta$, under
  the conditions arising in \cref{sec:window}; the development constructs the Euler system.
  \literatureform{special case}
\end{description}

\subsection{Representations of \texorpdfstring{$p$}{p}-adic, real and finite groups}\label{sec:literature:reps}

For a prime $g$ and $d\ge1$, a character $\theta$ of $\F_{g^d}^\times$ is \emph{regular} if
$\theta,\theta^g,\dots,\theta^{g^{d-1}}$ are pairwise distinct.

\begin{littheorem}[Green {\cite{Green1955}}]\label{thm:literature:green}
There is a map $\theta\mapsto\tau_\theta$ from characters of $\F_{g^d}^\times$ to functions on
$\GL_d(\F_g)$ such that: for regular $\theta$, $\tau_\theta$ is an irreducible cuspidal character of
degree $\prod_{i=1}^{d-1}(g^i-1)$, and $\tau_\theta(y)=(-1)^{d-1}\sum_{i=0}^{d-1}\theta(z^{g^i})$ for
every $y$ with irreducible characteristic polynomial and every root $z$ of it; for regular
$\theta,\theta'$, $\tau_\theta=\tau_{\theta'}$ if and only if $\theta'=\theta^{g^i}$ for some $i$; and every
irreducible cuspidal character of $\GL_d(\F_g)$ is some $\tau_\theta$.
\formal{GlobalTheory.greenRegularCharacters_holds}
\end{littheorem}

This is $\sfname{GreenRegularCharacters}$ \literatureform{same}.  It yields the proposition
$\sfname{Green}$ of \cref{def:green}, gives the irreducibility link (\cref{thm:linkIRR}), and fixes
the parameters of depth-zero supercuspidal representations in \cref{sec:galois}.

\begin{description}[leftmargin=0pt,itemsep=1pt,parsep=0pt,font=\normalfont\itshape]
\item[Bernstein--Zelevinsky theory and Zelevinsky's classification]
  \cite{BernsteinZelevinsky1976,BernsteinZelevinsky1977,Zelevinsky1980}.  For $\GL_n(F)$, $F$
  non-archimedean: admissibility, finite length of Jacquet modules and induced representations, the
  geometric lemma, irreducibility of $\pi_1\times\pi_2$ when no cuspidal representation in the
  support of $\pi_1$ is an unramified twist of one in the support of $\pi_2$; the Grothendieck ring
  is the polynomial ring on segments, every irreducible representation is $L(\mathfrak m)$ for a
  unique multisegment $\mathfrak m$, and standard modules have irreducible cosocle.
  \literatureform{same}
\item[Unitary representations] \cite{Tadic1986,Bernstein1984}.  Generic unitary representations
  of $\GL_N(F)$ are products of unitary square-integrable representations and complementary series;
  products of irreducible unitary representations are irreducible.  \literatureform{same}
\item[Depth zero] \cite{MoyPrasad1996,Borel1976,Casselman1980}.  A representation of $\GL_n(F)$ has
  depth zero if and only if its cuspidal support has; vectors fixed by $K_1(N\cap\GL_n(\OO_F))$,
  $K_1$ the first congruence subgroup, inject into the Jacquet module along $N$; Iwahori-fixed vectors
  and Jacquet modules (Borel--Casselman).  \literatureform{same}
\item[Gelfand--Kazhdan and Jacquet--Shalika]
  \cite{GelfandKazhdan1975,JacquetShalika1981a,JacquetShalika1981b}.  Over totally complex fields:
  $\pi^\vee\cong\pi\circ({}^t\cdot)^{-1}$ for cuspidal $\pi$ of $\GL_a$; strong multiplicity one for
  isobaric sums; $\Nm(w)^{-1/2}<|\alpha|<\Nm(w)^{1/2}$ for the Satake parameters of unitary cuspidal
  $\pi$.  \literatureform{special case}
\item[Local converse theorem] \cite{Henniart1993}.  Henniart's theorem, that cuspidal
  representations of $\GL_{m+1}(F)$ are determined by their $\varepsilon$-factors twisted by all
  irreducible representations of $\GL_m(F)$, is also proved ($\sfname{LocalConverseGL}$).
  \literatureform{same}
\item[Real unitary groups]
  \cite{HarishChandra1953,HarishChandra1966,HechtSchmid1975,Casselman1989,Wallach1992,DixmierMalliavin1978}.
  For $\Uni(p,q)$ and its finite covers: Harish-Chandra modules of irreducible unitary
  representations; discrete series for regular integral parameters $\lambda$, with lowest $K$-type
  $\lambda+\rho_n-\rho_c$, formal degree and exhaustion; holomorphic discrete series; Casselman--Wallach
  globalisation; Dixmier--Malliavin factorisation; Schur orthogonality.
  \literatureform{special case}
\item[Invariant theory and finite groups] \cite{Weyl1946,Howe1989a,Zelevinsky1981,Brauer1947}.
  Branching from $\Uni(n+1)$ to $\Uni(n)$; Howe duality for $(\GL_n(\C),\GL_{n'}(\C))$ on
  polynomials; restriction of cuspidal representations of $\GL_n(\F_q)$ to the mirabolic subgroup;
  Brauer's induction theorem (variant).  \literatureform{same unless indicated}
\end{description}

\subsection{Theta correspondence}\label{sec:literature:theta}

\begin{description}[leftmargin=0pt,itemsep=1pt,parsep=0pt,font=\normalfont\itshape]
\item[Weil representation] \cite{Weil1964,MoeglinVignerasWaldspurger1987,Kudla1994}.  The
  Stone--von Neumann theorem; the metaplectic representation of $\Sp_{2n}(\R)$; global Weil
  representations of unitary dual pairs with splitting characters; Kudla's splitting in the
  Schr\"odinger model for split $\Uni(\nu,\nu)$.  \literatureform{same}
\item[Siegel--Weil formula] \cite{Weil1965,KudlaRallis1988,Ichino2004}.  In Weil's range of
  absolute convergence: convergence of the theta integral and of the Siegel Eisenstein series, and
  their identity at the Siegel--Weil point ($\sfname{WeilSiegelFormulaUnitary}$,
  $\sfname{SiegelWeilUnitaryConvergent}$).  \literatureform{special case}  The regularised formula
  \cite{KudlaRallis1994} is not used: the comparison of theta integrals of complementary spaces needed
  in \cref{sec:families} is proved directly, by induction on the signature, from a Siegel--Weil
  formula for definite spaces.
\item[Howe duality] \cite{Howe1979,Howe1989}.  Archimedean Howe duality in the Fock model
  ($\sfname{HoweDualityReal}$); in the unramified case, vectors fixed by the hyperspecial subgroup
  of one member of the pair lie in the span of the translates of the vacuum vector by the other.
  General $p$-adic Howe duality \cite{GanTakeda2016} is not used.  \literatureform{special case}
\item[Rallis inner product formula] \cite{Rallis1984inj,GanQiuTakeda2014}.  In Weil's range of
  convergence, theta lifts of cusp forms whose lower lifts vanish are square-integrable, with Petersson
  products given by the doubling zeta integral ($\sfname{RallisInnerProductConvergent}$).
  \literatureform{variant}
\item[Doubling method] \cite{GelbartPiatetskiShapiroRallis1987,PiatetskiShapiroRallis1987,Langlands1976}.
  Continuation of Siegel Eisenstein series on $\Uni(V\oplus-V)$ with singularities of finite order,
  the basic identity and the continuation of the zeta integral (\literatureformin{variant},
  \cref{sec:literature:forms}); the unramified doubling integral for $\GL_n$,
  $L(s+\frac12,\alpha)L(s+\frac12,\alpha^{-1})\prod_{i=1}^n(1-q^{-2s-i})$ \literatureform{same}.
\end{description}

\subsection{Automorphic forms and base change}\label{sec:literature:aut}

\begin{description}[leftmargin=0pt,itemsep=1pt,parsep=0pt,font=\normalfont\itshape]
\item[Automorphic forms]
  \cite{Borel1969,HarishChandra1968,MoeglinWaldspurger1995,Flath1979,Clozel1990}.  Reduction theory
  for unitary groups; finite dimensionality of spaces of automorphic forms with fixed level,
  $K_\infty$-type and infinitesimal character; rapid decay of cusp forms; trace class of
  $R_{\mathrm{cusp}}(f)$; Flath's factorisation (variant); Clozel's purity lemma and cohomological
  criterion.  \literatureform{same unless indicated}
\item[Base change for unitary groups] \cite{Labesse2011}.  For a CM field $E$ with $E/E^+$
  unramified at all finite places, every automorphic representation of a definite $\Uni(n)$ has a
  weak base change to $\GL_n(\A_E)$, an isobaric sum of conjugate self-dual representations matching
  it at split places, at unramified inert places and at infinity
  ($\sfname{BaseChangeDefiniteUnitaryGroup}$).  The formal proof does not use the stabilised trace
  formula; it works with systems of Hecke eigenvalues at bounded level.  \literatureform{special case}
\item[Cyclic base change] \cite{ArthurClozel1989,Langlands1980}.  Base change and descent for
  $\GL_a$ along cyclic extensions of prime degree of totally complex fields; descent of $\GL_2$-forms
  from imaginary quadratic fields to $\Q$.  \literatureform{same}
\item[Hecke characters] \cite{Weil1956,Serre1968}.  The $\ell$-adic character attached to an
  algebraic Hecke character: existence, uniqueness, crystalline and de Rham properties with the
  expected Hodge--Tate weights, local--global compatibility (variant).
  \literatureform{same unless indicated}
\end{description}

\subsection{Shimura varieties and abelian varieties}\label{sec:literature:shimura}

\begin{description}[leftmargin=0pt,itemsep=1pt,parsep=0pt,font=\normalfont\itshape]
\item[Moduli spaces] \cite{MumfordFogartyKirwan1994,Kottwitz1992,HarrisTaylor2001,Zarhin1977}.
  Representability of Siegel moduli (polarisations whose kernel is killed by $d$, level $M\ge3$ odd
  and prime to $d$) by quasi-projective schemes over $\Z[\zeta_M,1/M]$; Kottwitz's integral models of unitary PEL Shimura varieties of the type of
  \cite{HarrisTaylor2001}: smooth quasi-projective fine moduli schemes of relative dimension $a(n-1)$
  over every Noetherian base ring carrying the base data; Zarhin's trick.
  \formal{GlobalTheory.UPel.HT.model_nonempty} \literatureform{variant}
\item[Semistable reduction] \cite{Grothendieck1972,SerreTate1968}.  An abelian variety over a
  $p$-adic field acquires semi-abelian reduction with split toric part after a finite extension;
  the N\'eron--Ogg--Shafarevich criterion.  \literatureform{same}
\item[Congruence relation] \cite{Wedhorn2000}.  A pointwise Eichler--Shimura congruence relation on
  Kuga--Sato varieties over integral models.  \literatureform{variant}
  The proof of \cref{thm:geoone} also uses a holomorphic Lefschetz formula for finite \'etale correspondences
  (\cref{thm:galois:hlf}). For the identity correspondence it is the formula of Atiyah and Bott
  \cite{AtiyahBott1968}; for correspondences we know no published proof (compare \cite{Tu2023}), so it is not
  listed here as a published result. It is proved in \cref{sec:pf:geo}.
\item[Toroidal compactifications]
  \cite{AshMumfordRapoportTai1975,FaltingsChai1990,Lan2013,Mumford1977,MoretBailly1985}.  Smooth
  projective toroidal compactifications of unitary Shimura varieties with normal crossings boundary
  and the logarithmic Kodaira--Spencer isomorphism ($\sfname{ToroidalForFanKS}$); integral toroidal
  compactifications of Siegel moduli over $\Z[\zeta_n,1/n]$ with the semi-abelian extension of the
  universal family and generated Hodge bundle ($\sfname{SiegelToroidalIntegral}$); Mumford's good
  metrics; extension of homomorphisms of semi-abelian schemes; subbundles of $\Lie G$ of a
  semi-abelian scheme over a curve have degree $\le0$.  \literatureform{same}
\item[Abelian schemes] \cite{Mumford1970,BerthelotBreenMessing1982,Messing1972}.  Over arbitrary
  rings, $H^1_{\mathrm{dR}}(A/R)$ is finite projective of rank $2g$ with Hodge filtration of rank
  $g$, compatible with base change, crystalline, with the duality pairing
  ($\sfname{AbelianDeRham}$); Grothendieck--Messing theory (\cref{thm:literature:gm}); dual abelian
  schemes (\cref{thm:literature:dual}); the theorem of the cube.
  \literatureform{same unless indicated}
\end{description}

\subsection{\texorpdfstring{$p$}{p}-adic Hodge theory and \'etale cohomology}\label{sec:literature:padic}

\begin{description}[leftmargin=0pt,itemsep=1pt,parsep=0pt,font=\normalfont\itshape]
\item[\'Etale cohomology] \cite{ArtinGrothendieckVerdier1972,Deligne1977}.  Proper base change,
  local acyclicity over absolutely integrally closed local domains (a case of smooth base change),
  finiteness for proper varieties, the projective bundle formula, specialisation isomorphisms.
  \literatureform{same; smooth base change: special case}
\item[Comparison theorems] \cite{Faltings1988,Faltings1989,FontaineMessing1987,Tsuji1999}.
  For smooth projective schemes over $\OO_K$, $K/\Q_\ell$ finite: the $\ell$-adic cohomology of the
  generic fibre is crystalline, with de Rham comparison and Frobenius compatibility; the
  Hodge--Tate decomposition.  \literatureform{special case}
\item[$p$-adic monodromy theorem] \cite{Berger2002,Fontaine1994}.  De Rham representations are
  potentially semistable ($\sfname{DeRhamPotSemistable}$); properties of $B_{\mathrm{cris}}$.
  \literatureform{same}
\item[Integral theory] \cite{BhattScholze2023,ScholzeWeinstein2013,AnschutzLeBras2023}.  Lattices in
  crystalline representations are realised by prismatic $F$-crystals, fully faithfully (stated
  through evaluation on the $A_{\inf}$-prism); prismatic Dieudonn\'e modules over $\OO_C$ with their
  Hodge filtration.  For lattices with small Hodge--Tate weights, Wach modules
  \cite{Wach1996,Berger2004} play the role of Fontaine--Laffaille theory \cite{FontaineLaffaille1982},
  in formulations adapted to the use.  \literatureform{variant}
\item[Local Bloch--Kato theory] \cite{BlochKato1990}.  The dimension formula for $H^1_f(K,V)$, $V$
  de Rham, and exact orthogonality of $H^1_f(K,V)$ and $H^1_f(K,V^*(1))$, rationally and integrally.
  \literatureform{same}
\end{description}

\subsection{Arithmetic and algebraic geometry}\label{sec:literature:geometry}

\begin{description}[leftmargin=0pt,itemsep=1pt,parsep=0pt,font=\normalfont\itshape]
\item[Arithmetic intersection theory] \cite{GilletSoule1990,BostGilletSoule1994,Philippon1991}.
  Arithmetic intersection numbers of hermitian line bundles and heights of cycles, with their
  calculus; Chow forms and their heights.  \literatureform{same}
\item[Arithmetic Hilbert--Samuel theorem and volumes]
  \cite{GilletSoule1992,AbbesBouche1995,Zhang1995,Moriwaki2009,Yuan2008}.  For a nef hermitian line
  bundle $\overline L$ on an arithmetic variety of dimension $D+1$ and any $\overline N$,
  $\log\#\{s\in H^0(L^m\otimes N):\|s\|_{\sup}\le1\}=\widehat{\deg}(\overline L^{D+1})\,m^{D+1}/(D+1)!+o(m^{D+1})$;
  $\widehat{\vol}(\overline L)=\widehat{\deg}(\overline L^{D+1})$ for nef $\overline L$, with $\ge$
  for relatively nef semipositive $\overline L$; Moriwaki's bigness criterion.
  \literatureform{variant}
\item[Resolution of singularities] \cite{Hironaka1964,BierstoneMilman1997,Wlodarczyk2005,Kollar2007}.
  Functorial resolution and embedded principalisation in characteristic $0$; in the form used, a
  reduced separated scheme of finite type over a field of characteristic $0$ is resolved by a
  sequence of blow-ups, isomorphic over the smooth locus, with simple normal crossings exceptional
  locus, and ideal sheaves on smooth varieties are principalised by blow-ups in smooth centres.
  \formal{FunctorialResolution.functorialResolutionCharZero} \literatureform{same}
\item[Vanishing and positivity]
  \cite{Kawamata1982,Viehweg1982,Nadel1990,AngehrnSiu1995,OhsawaTakegoshi1987,Lazarsfeld2004}.
  Kawamata--Viehweg vanishing, Kleiman's criterion and asymptotic Riemann--Roch for nef bundles
  \literatureform{same}; Nadel vanishing for pairs, Angehrn--Siu point separation, an $L^2$
  extension theorem of Ohsawa--Takegoshi type \literatureform{variant}; GAGA, Grothendieck's
  existence theorem and the Chern--Weil and Poincar\'e--Lelong formulas \cite{Serre1956,Grothendieck1961}
  \literatureform{same}.
\end{description}

\subsection{Statements in forms weaker or stronger than published}\label{sec:literature:forms}

We describe the formal statements which differ from the published theorems in ways that matter
for the proof or for reading the development.

\literaturepar{Local Langlands parameters}  The local Langlands correspondence for $\GL_n$
\cite{HarrisTaylor2001,Henniart2000,Scholze2013} is not used.  The proof uses instead the proposition
$\sfname{RecTExists}$ (\cref{def:rectexists,thm:rectexists}; compare \cref{rem:galois:llc}): at every
finite place $w$ of every number field
there is a map $\mathrm{rc}$ from irreducible smooth representations of the groups $\GL_n(E_w)$ to
Frobenius-semisimple Weil--Deligne representations which is local class field theory for $n=1$, is
compatible with the Langlands classification and with unramified twists, sends supercuspidal
representations to irreducible ones, and on \emph{good} representations (cuspidal support consisting
of characters or, when $E_w\cong\Q_g$, of twists of depth-zero supercuspidals) is compatible with
twists, central characters and contragredients, is injective, is tame in depth zero, and is given on
depth-zero supercuspidals by Green's parametrisation (\Cref{thm:literature:green}).  Bijectivity and
the preservation of $L$- and $\varepsilon$-factors of pairs are not required, and $\mathrm{rc}$ is
determined only on good representations.  Up to normalisations, $\sfname{RecTExists}$ follows from
the local Langlands correspondence together with its explicit description in depth zero
\cite{BushnellHenniart2005}, and it asserts much less; this is all that is used.  It is proved
directly from local class field theory, the theory of Bernstein--Zelevinsky and Zelevinsky,
Moy--Prasad and Green.  \formal{UpperTheory.Zero.LLC.recTExists}

\literaturepar{Galois representations}  In place of \cite[Thm.~2.1.1]{BarnetLambEtAl2014}
(\cref{thm:galois:blggt}), building on
\cite{HarrisTaylor2001,Shin2011,ChenevierHarris2013,Caraiani2012,Caraiani2014}, the proof uses the
proposition $\sfname{Row0b}$ (\cref{def:row0b,thm:row0b}): for a CM field $E$, every regular algebraic
cuspidal representation $\pi$ of $\GL_a(\A_E)$ polarised by a Hecke character of $E^+$ has a
continuous semisimple $r:G_E\to\GL_a(\Qlbar)$ with the Hodge--Tate weights given by the infinitesimal
character of $\pi$, such that $\mathrm{WD}(r|_{G_{E_v}})=\mathrm{rc}(\pi_v|\det|^{(1-a)/2})$ at places
$v\nmid\ell$ where the right side is abelian with $N=0$ (in particular at all unramified places),
the inertial characteristic polynomials agree at places $E_v\cong\Q_g$, $g>16a^2$, where it is tame
with square-free inertial characteristic polynomial and $N=0$, and $r$ is crystalline with the
expected weights at places above $\ell$ where $\pi$ is unramified (with a partial potentially
semistable compatibility at tame places), for every $\mathrm{rc}$ as in $\sfname{RecTExists}$.  The
polarisation of $r$, purity and local--global compatibility at general ramified places are not
asserted.  The development proves $\sfname{Row0b}$ through theta lifts and Kuga--Sato varieties
(\cref{sec:galois}), without the cited works.
\formal{UpperTheory.Gal.Int.galoisLocalGlobalStrong_recParameter}
Formulations of the full published theorems ($\sfname{LocalLanglandsGL}$,
$\sfname{GaloisRepOfPolarizedCuspidal}$) occur in the development but are not used; the structure
\texttt{OpenLiteratureStatements} taken by the main reduction consists of $\sfname{Row0b}$ and
$\sfname{RecTExists}$ and is instantiated by the two proofs cited.

\literaturepar{Deformations and duals of abelian schemes}  The following two statements are used in
the generality indicated.

\begin{theorem}[Deformations of abelian schemes]\Fmark\label{thm:literature:gm}
Let $R\to R_0$ be a surjection of commutative rings whose kernel carries nilpotent divided powers,
and consider abelian schemes over $R$ and their reductions to $R_0$.  For every such abelian scheme
the crystalline comparison map on its first de Rham cohomology is bijective and its Hodge filtration
is admissible; every admissible filtration is the Hodge filtration of a lift of the reduction; and a
homomorphism between reductions which respects the Hodge filtrations of two lifts lifts uniquely to a
homomorphism between these lifts.
\formal{GlobalTheory.Published.GrothendieckDeformation}
\end{theorem}

In \cite{Messing1972} a prime $p$ is assumed nilpotent on the base (the statement
$\sfname{GrothendieckMessing}$); \cref{thm:literature:gm} makes no such assumption.  For abelian
schemes the relevant crystal, defined through the universal vector extension, makes sense along
arbitrary nilpotent divided-power thickenings \cite{MazurMessing1974}.  For an abelian scheme $A\to S$
and an $S$-scheme $T$, call a line bundle $M$ on $A\times_ST$ \emph{additive} if
$(a+a',t)^*M\cong(a,t)^*M\otimes(a',t)^*M$ for all $S$-schemes $Y$ and $a,a':Y\to A$, $t:Y\to T$.

\begin{theorem}[Dual abelian schemes]\Fmark\label{thm:literature:dual}
Let $S$ be a scheme and $A\to S$ an abelian scheme of relative dimension $g$.  There are an abelian
scheme $B\to S$ of relative dimension $g$ and an additive line bundle $P$ on $A\times_SB$, trivial
on $A\times\{0\}$, such that for every $S$-scheme $T$ and every additive line bundle $M$ on
$A\times_ST$ there is a unique $x:T\to B$ with $M\cong(1\times x)^*P$; and the same holds with
$A$ and $B$ exchanged.
\formal{GlobalTheory.dualAbelianScheme}
\end{theorem}

The standard references construct the dual as a scheme when $A/S$ is projective
\cite{MumfordFogartyKirwan1994} or $S$ is normal \cite{FaltingsChai1990}, and as an algebraic space
in general; \cref{thm:literature:dual} ($\sfname{DualAbelianScheme}$) is stated and proved over an
arbitrary base scheme.  It enters \cref{sec:covolume} through Zarhin's trick and
\cref{thm:literature:gm}.

\literaturepar{Other differences}  (a) The doubling method is used without the statement, part of
\cite{GelbartPiatetskiShapiroRallis1987,PiatetskiShapiroRallis1987}, that the Siegel Eisenstein
series has at most a simple pole at $(b-n)/2$ when $n=\dim V<b$; accordingly, the field of
\texttt{UnitaryGroupThirdPlaceHypotheses} named after this simple pole has the true proposition as
its type.  (b) The Siegel--Weil and Rallis formulas
are used only in Weil's range of convergence; the archimedean theory only for $\Uni(p,q)$ and its
covers; global results on $\GL_a$ only over totally complex (in a few cases totally real) fields;
base change only when $E/E^+$ is unramified at all finite places.  (c) Kottwitz's models are
constructed over general Noetherian base rings, not only over the localised ring of integers of the
reflex field.  (d) Besides \cref{thm:literature:dual} and the field in~(a), two statements on theta
functions ($\sfname{ComplementaryThetaInert}$, $\sfname{ThetaRationalInvariance}$) have no precise
published counterpart.  (e) Some names in the development are historical:
$\sfname{ComplexMultiplicationJSt}$ is the main theorem of complex multiplication,
$\sfname{KStability}$ is homological stability for $\GL_n$, and several theorems whose names end in
\texttt{\_of\_literature} are new results of the development.

\subsection{Summary}\label{sec:literature:summary}

\Cref{tab:literature} lists the main results, the form of the formal statement and the sections in
which they are used.

{\footnotesize
\setlength{\tabcolsep}{4pt}
\setlength{\LTpre}{6pt}\setlength{\LTpost}{0pt}
\begin{longtable}{@{}>{\raggedright\arraybackslash}p{0.45\linewidth}>{\raggedright\arraybackslash}p{0.28\linewidth}>{\raggedright\arraybackslash}p{0.11\linewidth}>{\raggedright\arraybackslash}p{0.08\linewidth}@{}}
\caption{Main published results used in the proof; all of them are proved in the formal
development.}\label{tab:literature}\\
\toprule
Result & References & Form & Used in\\
\midrule
\endfirsthead
\toprule
Result & References & Form & Used in\\
\midrule
\endhead
\bottomrule
\endfoot
\multicolumn{4}{@{}l}{\textit{Analytic number theory}}\\*[1pt]
Functional equation of $L(\chi,s)$; $L(\chi,s)\ne0$ on $\Re s\ge1$ & \cite{Davenport2000,LoefflerStoll2025} & same & \S\ref{sec:reduction}\\*
Zero-free regions; Siegel; Siegel--Walfisz & \cite{delaValleePoussin1899,Siegel1935b,Walfisz1936} & same & \S\ref{sec:top}\\*
Linnik's theorem & \cite{Linnik1944} & same & \S\ref{sec:families}\\*
Hecke $L$-functions & \cite{Hecke1920,Tate1967thesis} & same & \S\S\ref{sec:families}--\ref{sec:window}\\
Siegel--Klingen; Deligne--Ribet & \cite{Siegel1970,DeligneRibet1980} & same; variant & \S\S\ref{sec:families}--\ref{sec:window}\\
Chowla--Selberg formula & \cite{ChowlaSelberg1967} & same & \S\ref{sec:covolume}\\
\literaturegroup{Class field theory and Galois cohomology}
Global and local class field theory & \cite{Artin1927,Tate1967gcft} & same & \S\S\ref{sec:families}, \ref{sec:galois}\\
Chebotarev's theorem (qualitative) & \cite{Chebotarev1926} & special case & \S\S\ref{sec:families}--\ref{sec:galois}\\
Complex multiplication & \cite{Shimura1971,Silverman1994} & same & \S\S\ref{sec:families}--\ref{sec:window}\\
Hermitian forms; strong approximation & \cite{Landherr1936,Kneser1966} & same & \S\S\ref{sec:families}, \ref{sec:covolume}\\
Local duality; Poitou--Tate; Euler characteristics & \cite{Tate1963,Milne2006} & same & \S\S\ref{sec:families}--\ref{sec:window}\\
Ferrero--Washington theorem & \cite{FerreroWashington1979} & same & \S\ref{sec:window}\\
Selmer groups of ordinary modular forms (Kato) & \cite{Kato2004} & special case & \S\ref{sec:window}\\
\literaturegroup{Representations}
Green's parametrisation of cuspidal characters & \cite{Green1955} & same & \S\S\ref{sec:top}, \ref{sec:families}, \ref{sec:galois}\\
Bernstein--Zelevinsky; Zelevinsky's classification & \cite{BernsteinZelevinsky1977,Zelevinsky1980} & same & \S\S\ref{sec:families}, \ref{sec:galois}\\
Generic unitary dual; unitary products & \cite{Tadic1986,Bernstein1984} & same & \S\ref{sec:families}\\
Depth zero and Jacquet modules & \cite{MoyPrasad1996} & same & \S\S\ref{sec:families}, \ref{sec:galois}\\
Gelfand--Kazhdan; Jacquet--Shalika & \cite{GelfandKazhdan1975,JacquetShalika1981a} & special case & \S\S\ref{sec:families}, \ref{sec:galois}\\
Discrete series of $\Uni(p,q)$; Casselman--Wallach & \cite{HarishChandra1966,Casselman1989} & special case & \S\ref{sec:families}\\
\literaturegroup{Theta correspondence and automorphic forms}
Weil representation; Kudla's splitting & \cite{Weil1964,Kudla1994} & same & \S\ref{sec:families}\\
Siegel--Weil formula (convergent range) & \cite{Weil1965,Ichino2004} & special case & \S\ref{sec:families}\\
Howe duality (archimedean, unramified) & \cite{Howe1979,Howe1989} & special case & \S\ref{sec:families}\\
Rallis inner product formula & \cite{Rallis1984inj,GanQiuTakeda2014} & variant & \S\ref{sec:families}\\
Doubling method; unramified doubling integral & \cite{GelbartPiatetskiShapiroRallis1987} & variant; same & \S\ref{sec:families}\\
Reduction theory; finiteness; rapid decay & \cite{Borel1969,HarishChandra1968} & same & \S\ref{sec:families}\\
Base change from definite $\Uni(n)$ to $\GL_n$ & \cite{Labesse2011} & special case & \S\S\ref{sec:families}--\ref{sec:window}\\
Cyclic base change and descent & \cite{ArthurClozel1989,Langlands1980} & same & \S\S\ref{sec:window}--\ref{sec:galois}\\
\literaturegroup{Shimura varieties and abelian varieties}
Siegel moduli; Kottwitz's integral models & \cite{MumfordFogartyKirwan1994,Kottwitz1992} & variant & \S\S\ref{sec:galois}--\ref{sec:covolume}\\
Semistable reduction & \cite{Grothendieck1972} & same & \S\S\ref{sec:galois}--\ref{sec:covolume}\\
Toroidal compactifications & \cite{AshMumfordRapoportTai1975,FaltingsChai1990,Lan2013} & same & \S\ref{sec:covolume}\\
Grothendieck--Messing theory & \cite{Messing1972,MazurMessing1974} & same; variant & \S\S\ref{sec:galois}--\ref{sec:covolume}\\
Dual abelian schemes & \cite{MumfordFogartyKirwan1994,FaltingsChai1990} & variant & \S\ref{sec:covolume}\\
\literaturegroup{$p$-adic Hodge theory and \'etale cohomology}
Base change and finiteness in \'etale cohomology & \cite{ArtinGrothendieckVerdier1972} & same & \S\ref{sec:galois}\\
Crystalline comparison; Hodge--Tate decomposition & \cite{Faltings1988,Faltings1989} & special case & \S\ref{sec:galois}\\
De Rham implies potentially semistable & \cite{Berger2002} & same & \S\S\ref{sec:families}, \ref{sec:galois}\\
Crystalline lattices and prismatic $F$-crystals & \cite{BhattScholze2023} & variant & \S\ref{sec:families}\\
Local Bloch--Kato theory & \cite{BlochKato1990} & same & \S\S\ref{sec:families}--\ref{sec:window}\\
\literaturegroup{Arithmetic and algebraic geometry}
Arithmetic intersection theory & \cite{GilletSoule1990,BostGilletSoule1994} & same & \S\ref{sec:covolume}\\
Arithmetic Hilbert--Samuel; volumes; bigness & \cite{GilletSoule1992,Moriwaki2009} & variant & \S\ref{sec:covolume}\\
Resolution and principalisation (characteristic $0$) & \cite{Hironaka1964,Kollar2007} & same & \S\ref{sec:covolume}\\
Kawamata--Viehweg vanishing; nef numerics & \cite{Kawamata1982,Viehweg1982} & same & \S\ref{sec:covolume}\\
Nadel vanishing; Angehrn--Siu & \cite{Nadel1990,AngehrnSiu1995} & variant & \S\ref{sec:covolume}\\
\end{longtable}
}

\endgroup
